\chapter{DOGMA reference model and falsifiable research program}

The previous four chapters introduced pieces: typed state, regulators, selective
transformation, structured memory, strands and reverse-complement actions. A list of
pieces is not yet a model. A model begins when every tensor has an owner, every
transition has a legal information set, every output has an objective, and a failed
experiment can force us to abandon a claimed benefit.

This chapter therefore freezes a small reference DOGMA. It is deliberately not the
future DOGMA Engine. It is a scientific object: simple enough to trace, differentiated
end to end, explicit about causality, and comparable with ordinary recurrent
baselines. The implementation is a proposal assembled from established ingredients
and earlier Evolutor contracts. Its value must come from experimentally useful
organization, not from biological vocabulary.
This freezes one additional reference variant. Chapter 18's exact pair-count
state and Chapter 20's block-clock bank remain separate teaching models:
learned loci here are not those objects renamed. Common contracts concern
ownership, legal information access, complete carry and falsifiable evaluation.
Architectural equivalence would require its own proof.

\section{The computational contract}

For batch size $B$, hidden width $D$, $K$ memory loci and sequence length $T$, let
$x_t\in\{0,\ldots,V-1\}^B$. The model-owned parameters include an embedding
$E\in\mathbb R^{V\times D}$ and learned affine maps. Request-owned state is
\[
h_t\in\mathbb R^{B\times D},\qquad
M_t\in\mathbb R^{B\times K\times D}.
\]
Both are zero-initialized at a request boundary unless a later runtime contract
explicitly supplies persistent state. A trace is observational metadata; it is not
silently fed back into the transition.

\Plate{01-contract-stack}{Reference DOGMA ownership and data flow. Biological
terms such as regulator and locus denote computational roles here, not literal
gene-expression chemistry.}{fig:stack}

With $e_t=E[x_t]$ and $q_t=[e_t;h_{t-1}]$, define a proposed fast-state update
\begin{align}
\tilde h_t &= \tanh(W_pq_t+b_p),\\
g_t &= \sigma(W_gq_t+b_g),\\
h_t &= (1-g_t)\odot h_{t-1}+g_t\odot\tilde h_t.
\label{evod22:fast}
\end{align}
This resembles established gated recurrence. The DOGMA claim is therefore
\emph{not} that Equation~\ref{evod22:fast} invents gating. Its role is to provide
a stable primitive on which the explicit memory and strand contracts can be tested.

Address the $K$ memory loci with
\[
a_t=\operatorname{softmax}(W_ah_t+b_a),\qquad
r_t=\sum_{k=1}^K a_{t,k}M_{t-1,k}.
\]
Writes use locus-wise gates $w_t=\sigma(W_wh_t+b_w)$ and value
$v_t=\tanh(W_vh_t+b_v)$:
\begin{equation}
M_{t,k}=(1-w_{t,k})M_{t-1,k}+w_{t,k}v_t.
\label{evod22:memory}
\end{equation}
The emitted feature is
$y_t=\tanh(W_m[h_t;r_t]+b_m)$. The reference code implements these equations
literally.

For one example use column vectors:
$W_p,W_g\in\mathbb R^{D\times2D}$,
$W_a,W_w\in\mathbb R^{K\times D}$,
$W_v\in\mathbb R^{D\times D}$ and
$W_m\in\mathbb R^{D\times2D}$. Bias lengths are $D,D,K,K,D,D$.
PyTorch batches examples as rows and applies transpose contractions inside
\texttt{nn.Linear}. A $C$-class head produces $C$ logits; causal base prediction
needs $C=4$, not the default two-class head.

\Listing{DOGMACell.step}

Reads use $M_{t-1}$, before the write at $t$. A new value cannot affect the
same step through memory. Reading $M_t$ instead changes the model even when
shapes agree. The read distribution addresses fixed locus indices, not
content-derived keys. Calling it content-addressed attention would erase a
meaningful distinction from Chapters 8 and 23.

\subsection{Trace a transition with numbers}

Set $D=K=1$, $g_t=1/2$, $\tilde h_t=1/2$, $w_t=1/4$ and $v_t=4/5$.
Choose $y_t=\tanh(h_t+r_t)$ and initialize both states to zero. Then
\begin{equation}
 h_t=\tfrac12(1-2^{-t}),\qquad
 M_t=\tfrac45(1-(3/4)^t),\qquad r_t=M_{t-1}.
 \label{evod22:trace}
\end{equation}
Substitution verifies the recurrence and initial conditions; induction proves
the formulas. Code fixes weights and biases to realize the constants and
generates this trace:
\begin{center}\input{\Generated/assigned_trace.tex}\end{center}
The first feature is $\tanh(1/4)$, not $\tanh(1/4+1/5)$. This independent
oracle detects read-before-write order, unlike copying the same update sequence
into both implementations.

\Plate{06-convex-write}{One memory-coordinate update: a write gate of one
quarter moves the old value $0.2$ toward the proposal $0.8$, producing $0.35$.
The updated memory belongs to the next carry, while this step's read uses
the old value. Positions on the line are values, not genomic coordinates.}{fig:write}

\subsection{Bounded state is not a gradient guarantee}

Starting from zero, each fast-state and memory coordinate remains in $[-1,1]$:
every update is a convex combination of its old coordinate and a bounded
$\tanh$ proposal. The addressed read is a convex combination of bounded rows.
This is a forward invariant, not a retention or gradient guarantee. Since
the gate and proposal depend on $h_{t-1}$, their derivatives add Jacobian
terms; the recurrent Jacobian is not simply $\operatorname{diag}(1-g_t)$.
Large learned weights can amplify perturbations despite bounded activations.

Including a biased $C$-class head, the model-owned parameter count is
\begin{equation}
 P=7D^2+(V+4)D+2K(D+1)+C(D+1).
 \label{evod22:parameters}
\end{equation}
Proposal/gate maps contribute $4D^2+2D$; addressing contributes $2K(D+1)$;
value/mixing maps contribute $3D^2+2D$; embedding and head supply the rest.
At $(V,D,K,C)=(4,12,3,2)$ this is $1208$. Carry stores $BD(K+1)$ floats;
histories, optimizer state, gradients, activations and input retention are
additional costs. Sharing parameters across strands does not remove the
second traversal's arithmetic.

\section{Mode one: causal streaming}

For causal prediction, the legal dependency set at position $t$ is the prefix
$x_{\le t}$ plus request state derived from that prefix. No reverse-complement
full-record feature is legal, because Chapter 21 showed that its aligned reverse
branch contains suffix information.

\Plate{02-causal}{Causal mode. Recurrent state moves only from earlier to later
positions. Appending or changing a suffix cannot change an already computed prefix
feature.}{fig:causal}

For next-token training, one legal objective is
\[
\mathcal L_{\rm LM}
=-\sum_{t=1}^{T-1}\log p_\theta(x_{t+1}\mid x_{\le t}).
\]
The target is shifted; the current token cannot be used as its own label. State is
reset between independent records unless the dataset explicitly defines continuity.

Chunking must carry \emph{both} $h$ and $M$. Method \texttt{run} returns a
complete \texttt{DOGMAState} without mutating caller-owned tensors. An empty
chunk emits a shape-$B\times0\times D$ history and preserves carry. Empty
offline records are rejected because mean pooling has no defined divisor.
This reference processes unpadded records; a later batching policy must define
whether invalid positions hold or advance state.

\Listing{DOGMA.run}

Consecutive chunk execution equals a single pass when it carries the complete
state, concatenates emitted features, and uses identical parameters and token
order. Induction on the token transition proves this. Tests compare outputs,
both carry tensors and parameter gradients at every split of a six-token record,
including empty end chunks. Detaching carry preserves forward values but truncates
gradients: that is a training policy, not a parity-preserving detail.
State-dependent gates also prevent a general replacement by Chapter 7's
input-only affine scan.

The test suite changes the final two tokens while preserving the first four and
checks exact prefix invariance to numerical tolerance. This is more informative than
checking only tensor shapes: it detects accidental suffix access.

\section{Mode two: offline strand-symmetric encoding}

Whole-record classification or annotation may legally use both orientations. Let
$R(X)$ reverse positions and complement A/C/G/T. Apply the \emph{same} causal
operator $F_\theta$ to $X$ and $R(X)$, realign the reverse branch, and combine:
\[
Z(X)=\big[F_\theta(X);\operatorname{rev}F_\theta(R(X))\big].
\]
The compact reference uses an invariant pooled feature
\[
\bar z(X)=\frac{1}{2T}\sum_t
\left(F_\theta(X)_t+
\operatorname{rev}F_\theta(R(X))_t\right).
\label{evod22:dual}
\]
Because parameters are shared and the two branches are exchanged by reverse
complement, the pooled feature is invariant up to floating-point arithmetic.

\Plate{03-offline-dual}{Offline dual mode shares parameters across both strand
orientations and realigns coordinates. The resulting full-record feature is
appropriate for offline objectives, not left-to-right generation.}{fig:dual}

This resolves the Chapter 21 choice rather than hiding it behind a flag. The two modes
have different legal objectives. If a task requires orientation-sensitive labels,
the output action must transform those labels rather than forcing invariance.

\section{Why call the memory structured?}

A single recurrent vector stores everything in one address. The reference adds
$K$ explicit loci and a learned address distribution. This does not by itself make
the system genomic. It creates a testable structural hypothesis: tasks with sparse,
reusable, differently timed information may benefit from explicit addressable loci
relative to a capacity-matched recurrent state.

\Plate{04-memory}{Fast state and explicit memory loci have different addresses but
share request ownership. Writes are gated; reads are weighted. Persistence across
requests is forbidden unless a runtime layer later supplies it explicitly.}{fig:memory}

The word \emph{locus} is an analogy with a precise endpoint. In biology a genomic
locus is a physical region with molecular regulation and inheritance. Here a locus
is one row of a floating-point tensor. The analogy motivates modular addressing;
it does not import chromatin, transcription, mutation or evolutionary semantics.

\section{Trace as evidence, not explanation}

For every step the reference can emit gate means and write-gate means.
The step numbers restart within each chunk. They are local diagnostics, not
a durable replay trace with record identity, absolute offsets and parameter hashes.
A richer
research implementation should record per-locus addresses, state norms, module
choices and output actions. But a trace is not automatically a causal explanation.
If a memory locus receives a high write gate whenever a motif appears, that is a
correlation. A stronger test intervenes: clamp the gate, permute loci, erase a locus,
or transplant state between matched examples and measure the output change.

A falsifiable mechanism claim therefore has three parts:
\[
\text{claimed role}\quad+\quad\text{intervention}\quad+\quad
\text{predicted change}.
\]
For example: ``locus 2 retains the first motif until the terminal decision'' predicts
that erasing locus 2 after the motif but before the decision will selectively reduce
accuracy on examples requiring that motif. If erasure has no effect across seeds,
the mechanistic story is wrong or incomplete.

\section{A small executable experiment}

The package constructs a synthetic whole-record task. A record is positive when it
contains motif ACG in either strand orientation; negatives are rejection-sampled so
neither orientation occurs. Positives are explicitly implanted. This task is chosen
because the offline dual mode has a legitimate symmetry to exploit. It is not a
proxy for regulatory genomics.

The integrated pilot trains DOGMA offline dual, a GRU, a flattened MLP and a
width-three CNN on the same 512 examples and evaluates on 128 held-out examples.
A local motif problem needs a local-convolution baseline and an exact string
oracle, not only weak generic comparisons. An independent ACG/CGT detector
checks every generated label and has perfect task accuracy by specification,
not through training.

Data seeds 137 and 139 are fixed before the model runs. Exact records and their
reverse complements are disjoint between splits. This excludes that synthetic
leakage channel but is not a homology-aware biological split. Initialization
seeds 11, 23 and 37 each run all model families for 30 full-batch epochs with
AdamW at learning rate $3\times10^{-3}$. Artifacts record every run, parameter
count, strand-pass count and the PyTorch environment.

\Listing{train_eval}
\begin{center}\input{\Generated/summary.tex}\end{center}
These are implementation-pilot means and sample standard deviations, not a
matched-budget ranking. Equal epochs do not equalize FLOPs, parameter capacity,
state storage or tuning opportunity. DOGMA traverses two orientations while
the learned comparison models traverse one. This table licenses no architecture
superiority claim; Chapter 25 supplies a shared comparison contract.

The generated \texttt{results.json} records actual losses, accuracies and parameter
counts from this package run. These results establish only that the reference
implementation can learn or fail on this controlled task under the recorded budget.
They do not establish an advantage on genomic data, language, long context or
production inference. The random-label control draws independent Bernoulli
targets for both training and testing, rather than permuting training labels
while scoring original motif labels. For independent uniform
$Y\in\{1,\ldots,V\}$ and predictor $q$,
\begin{equation}
 \mathbb E[-\log q_Y]=\log V+
 D_{\rm KL}(\operatorname{Uniform}(V)\Vert q)\ge\log V.
 \label{evod22:random}
\end{equation}
This is a population expectation conditional on an independent predictor,
not a bound on each finite-sample loss or training loss. Held-out values can
fall below $\log V$ by chance; training labels can be memorized. Audit
independence, uncertainty and repeated seeds before interpreting an exception.

A stronger study would add more seeds, tune each baseline under equal compute,
include a Transformer/SSM baseline where appropriate, stratify by motif position,
test longer sequences than training, and measure calibration as well as accuracy.

\section{Ablations that can kill the hypothesis}

The implemented read-path intervention replaces $r_t$ with zero while retaining
all parameter objects; the constructor requires $K\ge1$ and does not implement
a $K=0$ architecture. It preserves fast-state transitions and memory writes,
but changes emitted features after memory forms. Tests check these distinctions.
This is not a retrained, capacity-matched ablation. A second proposed ablation
freezes $w_t$ to a constant, testing whether input-dependent writes matter. The third
replaces the regulator $g_t$ with a learned constant vector. The fourth unties the
two strand branches in offline mode, testing whether symmetry is useful or merely
extra capacity.

These interventions must be interpreted with capacity controls. Removing memory
also removes parameters and computation; a weaker result may reflect smaller
capacity rather than the mechanism. Conversely, giving DOGMA two strand passes and
a baseline one pass is a compute advantage. Fair evidence reports these differences
or compensates for them.

\Plate{05-falsification}{A mechanism claim survives only if its benefit remains
after relevant baseline, budget and intervention controls. Negative evidence narrows
the architecture rather than being discarded.}{fig:falsify}

\section{Closest established mechanisms}

Gated recurrence is established. LSTM~\cite{lstm} introduced multiplicative gates to control
information flow through recurrent state; GRU-like models likewise use learned
update/reset behavior. Selective state-space models such as Mamba~\cite{mamba} make state updates
input dependent and design parallel algorithms around the resulting recurrence.
RWKV~\cite{rwkv} also explores the training/inference trade between parallel sequence processing
and recurrent execution.

DOGMA must therefore earn any distinction at a higher level of organization:
typed state ownership, explicit addressable loci, strand actions, trace/intervention
contracts, and eventually structural adaptation. If those pieces can be removed
without loss after fair capacity matching, the proposed architecture should collapse
toward the simpler established model. That is a scientifically useful outcome.

\section{Falsification matrix}

The research program can now be stated without slogans.

\begin{enumerate}
\item \textbf{Selective regulation.} Hypothesis: input-dependent gates improve tasks
with sparse state changes. Reject or narrow if constant-gate capacity-matched models
perform equivalently across seeds and interventions show no gate-specific effect.
\item \textbf{Addressable loci.} Hypothesis: explicit loci improve retention or
mechanistic locality. Reject if a same-budget recurrent state matches performance,
memory efficiency and intervention specificity.
\item \textbf{Strand sharing.} Hypothesis: parameter-tied dual processing improves
sample efficiency or consistency on strand-symmetric tasks. Reject if ordinary
augmentation achieves the same result at matched compute or if invariance hurts
orientation-sensitive tasks.
\item \textbf{Trace utility.} Hypothesis: typed traces support reproducible causal
diagnosis. Reject if trace variables are unstable across equivalent parameterizations
and interventions do not validate their claimed roles.
\item \textbf{Structural adaptation.} Deferred hypothesis: changing active modules
or topology can outperform fixed structure under a controlled adaptation budget.
Chapter 22 does not implement this and therefore makes no positive claim about it.
\end{enumerate}

\section{What the reference model does not contain}

There is no attention, no Transformer block, no DNA chemistry simulator, no
evolutionary search, no persistent agent memory, no distributed runtime and no
production kernel. Those omissions are intentional. Hermon DNA begins in Chapter 23
with a Transformer-based reference architecture. DOGMA Engine concerns begin in
Chapter 30. Evolutor's compiler/runtime sits above both much later.

The present chapter has one responsibility: make the non-Transformer proposal
specific enough that later engineering cannot retroactively change what an early
experiment meant.

\section{Exercises with worked answers}

\begin{enumerate}
\item Who owns $h_t$ and $M_t$? \textbf{Answer:} The request. Learned matrices and
embeddings are model-owned.
\item Why is zero initialization part of the scientific contract?
\textbf{Answer:} Hidden carry between independent records changes the function and
can leak data across examples.
\item Is Equation~\ref{evod22:fast} architecturally novel by itself?
\textbf{Answer:} No. Learned gated recurrence is established.
\item Why is offline dual state illegal for next-token prediction?
\textbf{Answer:} The aligned reverse branch contains suffix information.
\item Prove pooled Equation~\ref{evod22:dual} is reverse-complement invariant.
\textbf{Answer:} Applying $R$ swaps the two shared-parameter branches; their sum and
position average are unchanged after realignment.
\item Does invariance imply orientation-sensitive labels are impossible?
\textbf{Answer:} For an invariant pooled head, yes. Use an equivariant output action
or orientation metadata instead.
\item What does a high write gate prove?
\textbf{Answer:} Only that the parameterized model emitted a high gate. A causal
role requires intervention.
\item Why report parameter counts for baselines?
\textbf{Answer:} Apparent mechanism gains can be capacity gains.
\item What is the random-label control expected to show?
\textbf{Answer:} Independent held-out labels have chance population accuracy and
cross-entropy at least $\log V$. Finite samples fluctuate; memorizing training
noise establishes neither a violation nor the intended mechanism.
\item Design a locus-erasure intervention. \textbf{Answer:} Identify a time after
the candidate information is written, zero one locus for matched examples, continue
the same suffix, and compare output deltas against erasing other loci and sham edits.
\item When should DOGMA collapse toward a GRU/SSM?
\textbf{Answer:} When explicit loci, strand contracts and typed interventions add no
reproducible benefit after budget matching.
\item Design a genomic follow-up. \textbf{Rubric:} use a real provenance-tracked
dataset; split homologous/related sequences safely; define strand label action;
compare matched recurrent, SSM and Transformer baselines; use multiple seeds,
calibration and out-of-distribution lengths; pre-register ablations; distinguish
task accuracy from biological mechanism claims.
\item Derive parameter count for $D=12,K=3,V=4,C=2$.
\textbf{Answer:} $7(144)+8(12)+6(13)+2(13)=1208$. The orientations share
parameters but perform separate arithmetic.
\item Which parity fails when carry is detached between chunks?
\textbf{Answer:} Forward values can agree while gradients through prior writes
disappear. The full gradient comparison requires undetached complete carry.
\item Does an independent random-label test loss below $\log2$ prove leakage?
\textbf{Answer:} No. That is a population expectation bound. Check finite-sample
uncertainty, repeated seeds and independence before interpreting a deviation.
\item Compute the second feature of the assigned trace.
\textbf{Answer:} $h_2=3/8$ and $M_1=1/5$, so
$y_2=\tanh(0.575)\simeq0.519022$. Reading $M_2=0.35$ instead changes
transition order.
\end{enumerate}

\section{Vocabulary and next interface}

\Term{Request-owned state}: mutable state whose lifetime belongs to one inference or
training record unless persistence is explicitly declared. \Term{Model-owned
parameter}: learned state shared across requests. \Term{Mechanism intervention}:
a controlled change to an internal variable intended to test a causal role.
\Term{Falsification criterion}: an observation that would reject or narrow a claimed
benefit.

Chapter 23 starts a separate research direction inside Evolutor: Hermon DNA, a
Transformer-based DNA model. Keeping the two reference architectures explicit makes
their later comparison meaningful rather than turning ``DNA-native'' into a label
that can absorb any successful mechanism.
